% GENERATED FILE. Edit canonical lessons, then run npm run generate:books.
% canonical-source-hash: fnv1a64:e81f4a8fd91847a9
% canonical-lessons: RU-C357-sotrudnik, RU-C357-rukovoditel, RU-C357-podchinyonnyy, RU-C357-rukovodit, RU-C357-kollektiv

\chapter{Colleagues and managers}
\label{ch:ru-colleagues-and-managers}
\hlchaptermodality{357}

\begin{chapteropening}
\textbf{By the end of this chapter:} \emph{I can talk about the people I work with: colleagues, the manager, subordinates and the team, and say who is in charge.}

Complete the last lesson of chapter 357: I can talk about the people I work with: colleagues, the manager, subordinates and the team, and say who is in charge.
\end{chapteropening}

\section[\texorpdfstring{\ru{сотрудник} (sotrúdnik)}{сотрудник (sotrúdnik)}]{\ru{сотрудник} (sotrúdnik) — an employee, a member of staff}

\label{lesson:RU-C357-sotrudnik}



\begin{quote}
Before the new one: say the Russian for on the one hand, then the Russian for on the other hand.
\end{quote}

\subsection*{You'll want to know: \ru{сотрудник}}
\textbf{\ru{сотрудник}} (\emph{sotrúdnik}) — \textquotedblleft{}an employee, a member of staff\textquotedblright{}.

\begin{quote}
\ru{В} \ru{нашей} \ru{фирме} \ru{двадцать} \ru{сотрудников}.
\end{quote}

\begin{quote}
\emph{V náshey fírme dvádtsat sotrúdnikov.}
\end{quote}

\begin{quote}
There are twenty employees in our firm.
\end{quote}

Masculine; a woman is \emph{\ru{сотрудница}}. Formal; at work people also say \emph{\ru{коллега}}.

\begin{cousinweb}
\textbf{\ru{сотрудник}} is built from \emph{\ru{со}-} \textquotedblleft{}together\textquotedblright{} and \emph{\ru{труд}} \textquotedblleft{}work, labour\textquotedblright{}, the root of \emph{\ru{трудный}} \textquotedblleft{}hard\textquotedblright{}: one who works alongside you.
\end{cousinweb}

\begin{grammarlens}[title={What you've built}]
One more word for telling, explaining and describing at B1.
\end{grammarlens}

\subsection*{Your turn}
\begin{itemize}
\raggedright
  \item \emph{Say it:} \emph{sotrúdnik}
  \item \emph{Say it:} \emph{sotrúdnik}, once more
  \item \emph{Say it:} a sentence of your own with \emph{sotrúdnik}
  \item \emph{Recall:} say the Russian for on the one hand, then the Russian for on the other hand, then say \emph{sotrúdnik} again
\end{itemize}

\begin{tcolorbox}[breakable,colback=teal!4,colframe=teal!35!black,title={Before you move on}]
What does \ru{сотрудник} mean? (\textquotedblleft{}An employee, a member of staff\textquotedblright{}.) Say it once more.
\end{tcolorbox}
\section[\texorpdfstring{\ru{руководитель} (rukovodítel)}{руководитель (rukovodítel)}]{\ru{руководитель} (rukovodítel) — a manager, a head}

\label{lesson:RU-C357-rukovoditel}



\begin{quote}
Before the new one: say the Russian for on the other hand, then the Russian for an employee.
\end{quote}

\subsection*{You'll want to know: \ru{руководитель}}
\textbf{\ru{руководитель}} (\emph{rukovodítel}) — \textquotedblleft{}a manager, a head\textquotedblright{}.

\begin{quote}
\ru{Мой} \ru{руководитель} \ru{часто} \ru{ездит} \ru{в} \ru{командировки}.
\end{quote}

\begin{quote}
\emph{Moy rukovodítel chásto yézdit v komandiróvki.}
\end{quote}

\begin{quote}
My manager often goes on business trips.
\end{quote}

Masculine, also of a woman. \emph{\ru{Научный} \ru{руководитель}} — a thesis supervisor.

\begin{cousinweb}
\textbf{\ru{руководитель}} is built on \emph{\ru{руководить}}, a compound of \emph{\ru{рука}} \textquotedblleft{}hand\textquotedblright{} and \emph{\ru{водить}} \textquotedblleft{}to lead\textquotedblright{}, from Proto-Slavic \emph{*voditi} (EtymDB): one who leads others by the hand.
\end{cousinweb}

\begin{grammarlens}[title={What you've built}]
One more word for telling, explaining and describing at B1.
\end{grammarlens}

\subsection*{Your turn}
\begin{itemize}
\raggedright
  \item \emph{Say it:} \emph{rukovodítel}
  \item \emph{Say it:} \emph{rukovodítel}, once more
  \item \emph{Say it:} a sentence of your own with \emph{rukovodítel}
  \item \emph{Recall:} say the Russian for on the other hand, then the Russian for an employee, then say \emph{rukovodítel} again
\end{itemize}

\begin{tcolorbox}[breakable,colback=teal!4,colframe=teal!35!black,title={Before you move on}]
What does \ru{руководитель} mean? (\textquotedblleft{}A manager, a head\textquotedblright{}.) Say it once more.
\end{tcolorbox}
\section[\texorpdfstring{\ru{подчинённый} (podchinyónnyy)}{подчинённый (podchinyónnyy)}]{\ru{подчинённый} (podchinyónnyy) — a subordinate}

\label{lesson:RU-C357-podchinyonnyy}



\begin{quote}
Before the new one: say the Russian for an employee, then the Russian for a manager.
\end{quote}

\subsection*{You'll want to know: \ru{подчинённый}}
\textbf{\ru{подчинённый}} (\emph{podchinyónnyy}) — \textquotedblleft{}a subordinate\textquotedblright{}.

\begin{quote}
\ru{Начальник} \ru{добрый}, \ru{подчинённые} \ru{его} \ru{любят}.
\end{quote}

\begin{quote}
\emph{Nachálnik dóbryy, podchinyónnyye yegó lyúbyat.}
\end{quote}

\begin{quote}
The boss is kind; his staff love him.
\end{quote}

An adjective used as a noun: \emph{\ru{мой} \ru{подчинённый}, \ru{моя} \ru{подчинённая}}.

\begin{cousinweb}
\textbf{\ru{подчинённый}} is the participle of \emph{\ru{подчинить}} \textquotedblleft{}to put under someone’s authority\textquotedblright{}, built from \emph{\ru{под}} \textquotedblleft{}under\textquotedblright{} and \emph{\ru{чин}} \textquotedblleft{}rank, order\textquotedblright{}, the root of \emph{\ru{причина}} and \emph{\ru{чинить}}.
\end{cousinweb}

\begin{grammarlens}[title={What you've built}]
One more word for telling, explaining and describing at B1.
\end{grammarlens}

\subsection*{Your turn}
\begin{itemize}
\raggedright
  \item \emph{Say it:} \emph{podchinyónnyy}
  \item \emph{Say it:} \emph{podchinyónnyy}, once more
  \item \emph{Say it:} a sentence of your own with \emph{podchinyónnyy}
  \item \emph{Recall:} say the Russian for an employee, then the Russian for a manager, then say \emph{podchinyónnyy} again
\end{itemize}

\begin{tcolorbox}[breakable,colback=teal!4,colframe=teal!35!black,title={Before you move on}]
What does \ru{подчинённый} mean? (\textquotedblleft{}A subordinate\textquotedblright{}.) Say it once more.
\end{tcolorbox}
\section[\texorpdfstring{\ru{руководить} (rukovodít)}{руководить (rukovodít)}]{\ru{руководить} (rukovodít) — to manage, to be in charge of}

\label{lesson:RU-C357-rukovodit}



\begin{quote}
Before the new one: say the Russian for a manager, then the Russian for a subordinate.
\end{quote}

\subsection*{You'll want to know: \ru{руководить}}
\textbf{\ru{руководить}} (\emph{rukovodít}) — \textquotedblleft{}to manage, to be in charge of\textquotedblright{}.

\begin{quote}
\ru{Она} \ru{руководит} \ru{большой} \ru{командой}.
\end{quote}

\begin{quote}
\emph{Oná rukovodít bolshóy komándoy.}
\end{quote}

\begin{quote}
She manages a large team.
\end{quote}

Imperfective; present \emph{\ru{руковожу}, \ru{руководишь}}. It takes the instrumental: \emph{\ru{руководить} \ru{отделом}}.

\begin{cousinweb}
\textbf{\ru{руководить}} is a compound of \emph{\ru{рука}} \textquotedblleft{}hand\textquotedblright{} and \emph{\ru{водить}} \textquotedblleft{}to lead\textquotedblright{}, from Proto-Slavic \emph{*voditi} \textquotedblleft{}to lead\textquotedblright{} (EtymDB): to lead by the hand.
\end{cousinweb}

\begin{grammarlens}[title={What you've built}]
One more word for telling, explaining and describing at B1.
\end{grammarlens}

\subsection*{Your turn}
\begin{itemize}
\raggedright
  \item \emph{Say it:} \emph{rukovodít}
  \item \emph{Say it:} \emph{rukovodít}, once more
  \item \emph{Say it:} a sentence of your own with \emph{rukovodít}
  \item \emph{Recall:} say the Russian for a manager, then the Russian for a subordinate, then say \emph{rukovodít} again
\end{itemize}

\begin{tcolorbox}[breakable,colback=teal!4,colframe=teal!35!black,title={Before you move on}]
What does \ru{руководить} mean? (\textquotedblleft{}To manage, to be in charge of\textquotedblright{}.) Say it once more.
\end{tcolorbox}
\section[\texorpdfstring{\ru{коллектив} (kollektív)}{коллектив (kollektív)}]{\ru{коллектив} (kollektív) — a team, the staff}

\label{lesson:RU-C357-kollektiv}



\begin{quote}
Before the new one: say the Russian for a subordinate, then the Russian for to manage.
\end{quote}

\subsection*{You'll want to know: \ru{коллектив}}
\textbf{\ru{коллектив}} (\emph{kollektív}) — \textquotedblleft{}a team, the staff\textquotedblright{}.

\begin{quote}
\ru{У} \ru{нас} \ru{дружный} \ru{коллектив}.
\end{quote}

\begin{quote}
\emph{U nas drúzhnyy kollektív.}
\end{quote}

\begin{quote}
We have a close-knit team.
\end{quote}

Masculine. The people one works with, taken together.

\begin{cousinweb}
\textbf{\ru{коллектив}} is usually traced, through German, to Latin \emph{collēctīvus} \textquotedblleft{}gathered together\textquotedblright{}, from \emph{colligere} \textquotedblleft{}to gather together\textquotedblright{} (Lewis \& Short).
\end{cousinweb}

\begin{grammarlens}[title={What you've built}]
One more word for telling, explaining and describing at B1.
\end{grammarlens}

\subsection*{Your turn}
\begin{itemize}
\raggedright
  \item \emph{Say it:} \emph{kollektív}
  \item \emph{Say it:} \emph{kollektív}, once more
  \item \emph{Say it:} a sentence of your own with \emph{kollektív}
  \item \emph{Recall:} say the Russian for a subordinate, then the Russian for to manage, then say \emph{kollektív} again
\end{itemize}

\begin{tcolorbox}[breakable,colback=teal!4,colframe=teal!35!black,title={Before you move on}]
What does \ru{коллектив} mean? (\textquotedblleft{}A team, the staff\textquotedblright{}.) Say it once more.
\end{tcolorbox}
